\section{Process identification and the reflecting boundary}\label{app:boundary}

This appendix supplies the boundary facts used in Sections
\ref{sec:linear}--\ref{sec:returns}. The interior regularity theorem
is used only after a distributional equation has been established.
The reflecting condition is constructed separately.

\subsection{From the forward density to the reset process}

\begin{proof}[Proof of Proposition~\ref{prop:process}]
Let $F_t=S(t)h$. Put
$\Phi(v)=(\cos(2\pi v),\sin(2\pi v))$ and lift $F_t$ to a
probability $\nu_t$ on $\R^3$: choose $(v,y)$ with density $F_t$
and an independent fair sign, then set $(\xi,z)=(\Phi(v),\pm y)$.
This probability is supported on the unit circle times $\R$ and
need not have an ambient density. Away from the circle cut
$p=(1,0)$ use the inverse angle $v(\xi)\in(0,1)$ and define
\[
 B_\xi=2\pi J(v(\xi),|z|)(-\xi_2,\xi_1),\qquad
 B_z=m\,\operatorname{sign}z-z,\qquad
 A=\operatorname{diag}(0,0,2a).
\]
Choose any finite Borel angular drift at the cut and extend it
by zero off the circle. The value at the cut does not affect
the forward equation, since $\nu_t$ gives it zero mass.

For a compact smooth ambient test $\Psi$, its even pullback is
\[
 \varphi(v,y)=\tfrac12[\Psi(\Phi(v),y)+\Psi(\Phi(v),-y)].
\]
It has equal voltage endpoint values and $\varphi_y(v,0)=0$.
The reset terms cancel in the forward Green identity and the
low-conductance traces vanish. Integration by parts in $y$ has
no remaining floor contribution. The result is precisely the
Euclidean Fokker--Planck equation with drift $B$ and covariance
$A$, whose second-order term is $(1/2)A:D^2=a\partial_{zz}$.
The lifted probabilities are narrowly continuous and
\[
 \int_0^t\int_{\R^3}(|B|+|A|)\dd\nu_s\dd s<\infty,
\]
since $|B|\le C(1+|z|)$ on their support and $F$ has a bounded
first conductance moment.

Trevisan's superposition theorem \cite[Definition 2.2,
condition (4), and Theorem 2.5]{Trevisan15} supplies a continuous
martingale solution with these one-time laws. The angular
component lies on the circle at all rational times, and hence
at all times by continuity. Localized coordinate and product
tests show that its angular martingale has zero quadratic
variation, while its signed conductance satisfies
\begin{equation}\label{eq:signed-conductance}
 \dd Z_t=(m\,\operatorname{sign}Z_t-Z_t)\dd t+\sqrt{2a}\dd W_t .
\end{equation}
Here $W$ is scalar Brownian motion in the full solution
filtration, by the martingale quadratic variation and the scalar
L\'evy characterization. The stochastic integral, maximal
inequality, localized It\^o formula, and scalar characterization
used here can, for example, be taken from
\cite[Theorems 2.32, 4.10, 4.17, 4.19--4.20 and the scalar case
of Theorem 5.11]{DeVecchi22}.

We justify folding directly. Stop at
$\tau_R=\inf\{t:|Z_t|\ge R\}$ and apply the localized It\^o formula
to $\varphi_\delta(z)=\sqrt{z^2+\delta^2}$, made bounded outside
$[-R-1,R+1]$. For $0<\delta<1$,
\[
 a\E\int_0^{t\wedge\tau_R}\varphi_\delta''(Z_s)\dd s
       \le2R+2+(m+R)t .
\]
On $|z|\le\delta$ its second derivative is at least
$(2\sqrt2\,\delta)^{-1}$. Hence
\[
 \E\int_0^{t\wedge\tau_R}\one_{\{|Z_s|\le\delta\}}\dd s
       \le C_{a,R,t}\delta .
\]
Let $\delta\downarrow0$ and then $R\uparrow\infty$ to conclude
that $Z$ spends zero time at zero. The drift terms in the
same It\^o formula converge to those for $|Z|$; the stochastic
integrals converge uniformly in probability by isometry and
the maximal inequality. Consequently
\[
 K_t=|Z_t|-|Z_0|-\int_0^t(m-|Z_s|)\dd s
       -\sqrt{2a}\int_0^t\operatorname{sign}(Z_s)\dd W_s
\]
is a continuous nondecreasing limit of
$a\int_0^t\varphi_\delta''(Z_s)\dd s$.
On a path interval where $|Z|\ge d>0$, the latter increments
are bounded by $at\delta^2/d^3$ and vanish. Thus $K$ increases
only when $|Z|=0$. The martingale
$W'_t=\int_0^t\operatorname{sign}(Z_s)\dd W_s$ has quadratic
variation $t$, so it is scalar Brownian motion in the same
filtration. The folded $Y=|Z|$ satisfies
\eqref{eq:reset-process}.

The half-line Skorokhod map
$x(t)\mapsto x(t)+\sup_{s\le t}(-x(s))_+$
has supremum-norm Lipschitz constant at most two. Its integral
equation with drift $m-Y$ is a contraction on any interval of
length less than $1/2$. Iteration proves pathwise uniqueness
and makes $Y$ a measurable function of $Y_0,W'$. Brownian
increments are independent of the initial full sigma-field.
If $p(y)=\int_0^1h(v,y)\dd v$, Fubini therefore shows that the
conditional density of $V_0$ given the entire future $Y$ path
is $h(v,Y_0)/p(Y_0)$ on the set $p(Y_0)>0$.

For a fixed continuous conductance path, before a reset,
\begin{equation}\label{eq:voltage-flow}
 V_t=A_tV_0+C_t,\quad
 A_t=\exp\!\left[-\int_0^t(Y_s+1)\dd s\right],\quad
 C_t=\int_0^t\exp\!\left[-\int_s^t(Y_r+1)\dd r\right]2Y_s\dd s .
\end{equation}
The $C^1$ function $H(t)=(1-C_t)/A_t$ satisfies
$H'(t)=-J_F(Y_t)/A_t$. A grazing threshold hit would make
$V_0$ a critical value of $H$. Critical values of a $C^1$
real function on a bounded interval have Lebesgue measure
zero: cover the zero-derivative set by open intervals where
$|H'|<\epsilon$ and bound the lengths of their images by
$\epsilon$ times their total length, then let $\epsilon$
decrease to zero. Exhaust the time intervals.
The conditional density just proved excludes such a hit.
A first hit from below has $J_F\ge0$, and thus in fact $Y>1$.

After a reset the initial voltage is zero, so a separate
point-start argument is required. Until a floor encounter the
auxiliary unreflected conductance and voltage flow have a joint
density at every positive time. Indeed It\^o's formula gives
their distributional forward equation. The smooth fields
$\partial_y$ and
$\partial_t+J\partial_v+(m-y)\partial_y$, and their bracket
$(2-v)\partial_v-\partial_y$, span all time--space directions
for $v<2$. H\"ormander's theorem
\cite[Theorem 1.1]{Hormander67} gives a smooth space--time
density there. The explicit flow stays below $2$; narrow
continuity against compact tests upgrades the almost-every-time
density identification to every positive time.

Starting after a reset with $Y>1$, wait $\delta=1/n$ before the
next hit or first floor encounter. The law at $\delta$ restricted
to this survival event is dominated by the auxiliary law and
has a joint density. Future Brownian increments are independent
of the past event. The conditional-density critical-value
argument therefore applies after $\delta$.
The countable collection of these events covers all possible
finite next hits. Independence after stopping times follows by
approximating them from above by dyadic stopping times and
using Brownian continuity and deterministic-time independence.
Thus the argument repeats after every reset.

Before its first cut encounter the circle path inverts to the
ordinary voltage equation. It cannot reach zero from inside,
since $V_t\ge V_0A_t>0$. At a non-grazing outgoing hit the
angular speeds on both sides of the cut are positive, as
$J_F>0$ and $J_0=2Y>0$. Continuity of $Y$ makes them bounded
away from zero in a neighborhood. The path crosses in the
positive direction and enters the unique voltage trajectory
from zero. It cannot stick on an interval; in addition its
total occupation of the cut is zero by the one-time marginals.
The pathwise identity for $X$ gives \eqref{eq:count-bound}.
Since $Y$ is continuous, there is no finite-time accumulation
of hits. The same bound and the first conductance moment make
the expected count finite. This identifies the circle process
with the reset process, proving its law assertion.

For smooth $q$ supported in $(0,t)\times(1,\infty)$ apply
piecewise It\^o to $\varphi(s,v,y)=vq(s,y)$, whose reset jump
is $q$. Comparing with the forward Green formula gives
\eqref{eq:count-measure} against these tests. The trace of a
nonnegative forward density is nonnegative in its kinetic
trace measure; this also follows by the bounded negative-part
renormalization in the same Green identity. Both measures are
therefore locally finite positive measures. All hits have
$Y>1$. A spike at a deterministic time would put the circle
path at the cut, which its marginal excludes. There is no
initial spike for a density. Exhaust the conductance and
time endpoints. Finite expected count controls the total
measure, and monotone convergence extends equality first to
bounded Borel tests and then to integrable tests. This proves
the complete statement.
\end{proof}

\subsection{Local backward equations, including the floor}

We give the local construction for killed values with the weights
of Section~\ref{sec:returns}. Their interior sources belong to the
smooth class specified before Lemma~\ref{lem:cycle-energy} and obey
the stated Gaussian weighted bound. We also treat bounded
full-process resolvents with compact smooth source.
In both cases write
\begin{equation}\label{eq:backward-source-equation}
 \G u=\lambda u-r .
\end{equation}
The source is zero for a killed harmonic. Nonzero killed sources
used below are the constant source in the static block and the
smooth corrected-mode residual. The full-process case has
$\operatorname{Re}\lambda>0$ and compact smooth bounded $r$.

First establish the interior distributional equation. On an interior
rectangle use the unreflected, unreset process until exit from a larger
rectangle. Its conductance is explicit and its voltage is
\eqref{eq:voltage-flow}; the stochastic flow is triangular in the
initial variables, with positive Jacobian $e^{-t}A_t$. For a compact
smooth initial density $\varphi$, Taylor expansion of its inverse flow
and inverse Jacobian gives
\begin{equation}\label{eq:local-transpose-expansion}
 P_h^T\varphi
  =\varphi+h[-\partial_v(J\varphi)
             -\partial_y((m-y)\varphi)+a\varphi_{yy}]+o(h)
\end{equation}
in local $L^1$ and in the fixed Gaussian-weighted $L^1$ used to pair
with $u$. To see the remainder explicitly, the conductance increment
has mean $(m-y)h+O(h^2)$, variance $2ah+O(h^2)$, and $j$th absolute
moment $O(h^{j/2})$; its random contribution to the voltage increment
is $O(h^{3/2})$ in every fixed finite moment. The inverse Jacobian
is $1+h(y+2)+O_{L^p}(h^{3/2})$. Bounded third derivatives of
$\varphi$ control the Taylor remainder. Outside an enlarged compact
set the transported density has a Gaussian small-time tail, which
dominates the fixed exponential-quadratic weight. The probability
of leaving the larger rectangle before $h$ is $O(e^{-c/h})$;
higher stopped weights make its contribution $o(h)$.
For a smooth killed source the reward increment is $hr+o(h)$ on
compact sets of starts. The same localization and the higher-weight
reward bounds justify this term in the paired limit.
Dynamic programming, pairing with $\varphi$ and division by $h$,
now give \eqref{eq:backward-source-equation} in distributions.

We use the real-coefficient regularity theorem literally at complex
frequency. Introduce a real variable $z$ and set
$U(z,v,y)=e^{-\lambda z}u(v,y)$. Then
$(\partial_z+\G)U=-e^{-\lambda z}r$.
The fields $\partial_y$ and
$\partial_z+J\partial_v+(m-y)\partial_y$, with the same bracket
$(2-v)\partial_v-\partial_y$, span all three directions.
Apply \cite[Theorem 1.1]{Hormander67} to the real and imaginary
parts of $U$ separately. Thus $u$ is interior smooth.

At $v=0,y>0$ extend the auxiliary voltage a fixed small distance
below zero. An interior physical path never exits through zero;
the high-conductance block argument still bounds the hitting
time from the added finite voltage interval. The weak equation
and interior regularity on this extension make $u(0,y)$ smooth
for $y>0$. For a constant reset jump, the outgoing datum is
therefore smooth on $y>1$. Its positive outgoing speed gives
the terminal trace. At $v=1,y<1$, choose a small rectangle
straddling that face on which $J\le-j_0<0$. Stop the auxiliary
process on its exit. Its exit time is bounded by the voltage
width divided by $j_0$. Retain the original data on reachable
physical sides and extend side data continuously to the added
part. A physical start cannot exit through the added upper
voltage side. The local value agrees with $u$ inside, solves
the weak equation across the face, and is smooth there.
The deterministic exit bound suffices for any fixed complex
discount; no global hitting estimate outside the physical
voltage interval is assumed.

Point starts at positive conductance have no grazing hit by the
positive-time density argument above. For a start at $y=0$,
before any finite threshold hit the conductance first reaches
$1/2$ strictly earlier; start the same argument at that stop.
Continuous dependence of the Skorokhod and voltage equations
then implies almost-sure continuity of $\tf,Y_\tf$ in the
initial state at a transverse hit. Before the hit the voltage
on each shorter compact interval stays a positive distance
below threshold, and the final transverse crossing persists.
The stopped $4\beta$ weights and stronger discount moment give
uniform integrability on compact sets of starts. Thus killed
values are continuous, including at the floor. For a bounded
full-process source, the same coupling on finite intervals
and the discounted tail prove continuity.

Fix $0<v_-<v_+<1$ and choose $y_c>0$ so that $q=-J$ satisfies
$0<q_-\le q\le q_+$ on the resulting floor rectangle.
The backward equation becomes the forward voltage-clock problem
\begin{equation}\label{eq:floor-clock}
 q u_v=au_{yy}+(m-y)u_y-\lambda u+r,\qquad
 u_y(v,0)=0 .
\end{equation}
Give it the continuous initial data at $v_-$ and the upper-side
data at $y_c$. After a smooth lifting of that side, use a mixed
Neumann--Dirichlet basis on $(0,y_c)$. Its $q$-weighted mass
matrix is uniformly positive. Testing the finite systems by
their solutions yields
\[
 \tfrac12\partial_v\int_0^{y_c}q|w|^2\dd y
        +a\int_0^{y_c}|w_y|^2\dd y
 \le C\int_0^{y_c}(|w|^2+|\text{source}|^2)\dd y .
\]
The drift floor term is $-m|w(v,0)|^2/2$, with favorable sign.
Weak Hilbert limits, the variational equation, and the
zero-data energy inequality give existence and uniqueness,
with continuous $L^2$ time representative. Testing with $w_v$
after a positive voltage interval bounds $w_v,w_{yy}$ in $L^2$;
repeated differentiated estimates give the local boundary
regularity needed for It\^o. The absence of a diffusion floor
term in the variational identity encodes its Neumann trace.
For continuous initial data, approximation and the maximum
bound give the same solution. The full Galerkin passage and
positive-time arguments are detailed for the analogous
Dirichlet clock in Appendix~\ref{app:clock}; here the variational
space imposes only the upper-side zero value.

Until this rectangle is exited, physical voltage decreases
strictly. In its inverse clock the conductance generator is
$q^{-1}[a\partial_{yy}+(m-y)\partial_y]$, its discount is
$\lambda/q$, and its source is $r/q$. It\^o applied to the
constructed solution cancels the drift and the Neumann
derivative cancels the regulator. Physical exit time is
bounded by $(v_+-v_-)/q_-$. Stopped dynamic programming has
the same boundary data, so this solution equals $u$.
This proves the actual floor derivative for $v>0$, without
applying an interior theorem at a boundary.

\subsection{Growth for a constant jump and a bounded-source energy}

We first supply the growth assertion used in
Proposition~\ref{prop:constant-jump}. Suppose the reset jump is $A$.
Choose $p>\max\{1,r_0\}$ and a fixed $y_c>m,1$ so large that
\[
 (\G_y+r_0)y^p\le-d(y+r_0L),\qquad
 \G_y=(m-y)\partial_y+a\partial_{yy},\quad y\ge y_c .
\]
This follows from
$\G_y y^p=-py^p+pmy^{p-1}+ap(p-1)y^{p-2}$.
For $S=\inf\{t:Y_t\le y_c\}$, stopped It\^o gives
\[
 \E e^{r_0S}+
 \E\int_0^S e^{r_0t}(Y_t+r_0L)\dd t\le C(1+y^p).
\]
Recall $X=\log[2/(2-V)]\in[0,L]$, $L=\log2$.
Between hits $\dot X=Y-(e^X-1)$, and each reset decreases $X$ by $L$.
Including the reset jumps in $e^{r_0t}X_t$ gives
\[
 L\sum_{t_j\le S}e^{r_0t_j}
       \le X_0+\int_0^S e^{r_0t}(Y_t+r_0L)\dd t .
\]
Patching the killed martingales until $S$ consequently yields
\[
 u(z)=\E_z[e^{-\lambda S}u(Z_S)]
                +A\E_z\sum_{t_j\le S}e^{-\lambda t_j}.
\]
To remove upper-conductance and time localizations before asserting
this equality, choose $\beta<\beta'<1/2$ and enlarge $y_c$ so that
$(\G_y+r_0+d_0)w_{\beta'}\le0$ above it. The upper-exit term
is bounded by a fixed starting-state constant times
$w_\beta((H-m)/\varepsilon)/w_{\beta'}((H-m)/\varepsilon)$,
which tends to zero. The time endpoint is bounded by
$e^{-d_0t}w_{\beta'}(\eta_0)$.
The displayed count bound controls the spike sum independently.
Since $u$ is bounded on $y=y_c$, this proves
\begin{equation}\label{eq:polynomial-growth}
 |u(v,y)|\le C_{a,m,\lambda}(\norm f_\beta+|A|)(1+y^p).
\end{equation}
A bounded interior source adds a controlled
$\E\int_0^S e^{r_0t}|r|\dd t$ and gives the same conclusion.
These are fixed-$a$ growth constants.

The independent dual estimate of Section~\ref{sec:linear} needs
only bounded full-process source values. We state it separately
to keep its proof independent of killed-value gradients.

\begin{proposition}[Bounded-source backward energy]\label{prop:backward-source}
For bounded compact smooth $r$ and $\lambda=s+i\theta$ with $s>0$,
let
\[
 u(z)=\E_z\int_0^\infty e^{-\lambda t}r(Z_t)\dd t .
\]
It is bounded by $\norm r_\infty/s$, has zero reset jump,
satisfies \eqref{eq:backward-source-equation}, and has the local
Neumann floor constructed above. For $s$ larger than a fixed
constant,
\begin{equation}\label{eq:backward-energy}
 (s-C)\norm u_{\HH_{-\sigma}}^2+
          a\norm{u_y}_{\HH_{-\sigma}}^2
 \le \norm r_{\HH_{-\sigma}}\norm u_{\HH_{-\sigma}},
\end{equation}
uniformly in $\theta$. For $s$ also larger than the forward
growth bound, this value is the transpose of the full forward
resolvent.
\end{proposition}

\begin{proof}
The local construction above applies to this bounded value,
including its homogeneous reset and actual floor derivative.
Let $\rho=e^{-2\sigma y}$. On an exhausted domain the real
energy identity has volume coefficient
\[
 C_\sigma(y)=(\tfrac12-\sigma)y+1+\sigma m+2a\sigma^2,
\]
gradient term $-a\int|u_y|^2\rho$, floor term
$-(m/2+a\sigma)\int|u(v,0)|^2\dd v$, and voltage term
$\frac12\int\rho[J_F|u_F|^2-J_0|u_0|^2]\dd y$.
The floor and low-conductance voltage terms are nonpositive.
At high conductance $u_F=u_0$ and $J_F<J_0$, so that term is
nonpositive too. Since $\sigma>1$, $C_\sigma$ is bounded above.

Here is a localization that proves the first global gradient
bound rather than assuming it. Remove rectangles of voltage
width $\delta^3$ and conductance width $\delta$ around $(1,1)$
and $(0,0)$. Choose cutoffs $\chi_\delta$ with
$|\partial_v\chi_\delta|\le C\delta^{-3}$,
$|\partial_y\chi_\delta|\le C\delta^{-1}$ and zero floor
derivative. Add a large-$y$ cutoff $\zeta_R$ and test with
$\chi_\delta^2\zeta_R^2\overline u\rho$.
Near a corner $u$ is bounded, $|J|=O(\delta)$, and area is
$O(\delta^4)$. The voltage-cutoff error is $O(\delta^2)$,
the squared conductance-cutoff error is $O(\delta^2)$, and the
drift-cutoff error is $O(\delta^3)$. Young's inequality absorbs
the unknown gradient cross terms before any global gradient
integrability is assumed. A reset-cutoff mismatch is confined
to an outgoing interval with $\int J_F\dd y=O(\delta^2)$.
At large $y$, boundedness of $u$ and the exponential weight
remove all tail errors after the same absorption.

The resulting gradient bounds are uniform in $\delta,R$,
so Fatou gives $u_y\in\HH_{-\sigma}$. Repeat the passage with
this known gradient: the first-derivative cutoff errors now
tend to zero by Cauchy--Schwarz, since
$\norm{(\chi_\delta)_yu}_2=O(\delta)$.
The exact energy inequality and the preceding signs yield
\eqref{eq:backward-energy}. Finally Proposition~\ref{prop:process}
and Fubini identify the value paired against any positive
initial density with the full forward Laplace resolvent.
Splitting signed data and density in the weighted spaces
give the claimed transpose identification.
\end{proof}
